\documentclass[10pt,a4paper]{article}
\usepackage[utf8]{inputenc}
\usepackage[T1]{fontenc}
\usepackage[ngerman]{babel}
\usepackage{geometry}
\usepackage{xcolor}
\usepackage{tabularx}
\usepackage{array}
\usepackage{booktabs}
\usepackage{longtable}
\usepackage{colortbl}
\usepackage{enumitem}
\usepackage{helvet}
\usepackage{tikz}
\usepackage{tcolorbox}
\usepackage{fancyhdr}
\usepackage{titlesec}
\usepackage{hyperref}
\usepackage{pifont}
\usepackage{multirow}

\usetikzlibrary{shapes.geometric, arrows.meta, positioning}

\renewcommand{\familydefault}{\sfdefault}
\geometry{margin=2cm, top=2.5cm, bottom=2.5cm}
\setlength{\parindent}{0pt}
\setlength{\parskip}{6pt}
\setlength{\headheight}{14pt}

% Colors
\definecolor{spanisch}{HTML}{c41e3a}
\definecolor{lightred}{HTML}{fdf2f4}
\definecolor{darkgray}{HTML}{333333}
\definecolor{lightgray}{HTML}{f5f5f5}
\definecolor{spaet}{HTML}{8e44ad}
\definecolor{lightspaet}{HTML}{f5eef8}
\definecolor{success}{HTML}{27ae60}
\definecolor{warning}{HTML}{f39c12}
\definecolor{linkblue}{HTML}{2c5aa0}
\definecolor{englishgray}{HTML}{666666}

% ============================================================================
% BILINGUAL MACROS / Zweisprachige Makros
% ============================================================================
% \bitext{Deutscher Text}{English text}
\newcommand{\bitext}[2]{%
    #1\\[2pt]%
    {\small\color{englishgray}\textit{#2}}%
}

% \bipar{Deutscher Absatz}{English paragraph}
\newcommand{\bipar}[2]{%
    #1\par\vspace{2pt}%
    {\small\color{englishgray}\textit{#2}}\par\vspace{6pt}%
}

% Headers
\pagestyle{fancy}
\fancyhf{}
\fancyhead[L]{\small\textcolor{spaet}{\textbf{Spanisch Spätbeginner MV}} {\scriptsize\color{englishgray}| \textit{Spanish Late Beginners}}}
\fancyhead[R]{\small\textcolor{darkgray}{v2.0 | Januar 2026}}
\fancyfoot[C]{\thepage}

% Section formatting
\titleformat{\section}{\Large\bfseries\color{spaet}}{}{0pt}{}[\vspace{-4pt}\textcolor{spaet}{\rule{\textwidth}{1.5pt}}]
\titleformat{\subsection}{\large\bfseries\color{darkgray}}{}{0pt}{}
\titleformat{\subsubsection}{\normalsize\bfseries\color{spanisch}}{}{0pt}{}

% Hyperref
\hypersetup{
    colorlinks=true,
    linkcolor=spaet,
    urlcolor=linkblue
}

% Custom boxes
\tcbset{
    colback=lightgray,
    colframe=spaet,
    boxrule=1pt,
    arc=2pt,
    left=6pt, right=6pt, top=4pt, bottom=4pt
}

\newtcolorbox{wichtigbox}{colback=lightred, colframe=spanisch, fonttitle=\bfseries, title=WICHTIG / \textit{IMPORTANT}, toptitle=3pt, bottomtitle=3pt}
\newtcolorbox{empfehlungbox}{colback=lightspaet, colframe=spaet, fonttitle=\bfseries, title=HANDLUNGSEMPFEHLUNG / \textit{RECOMMENDATION}, toptitle=3pt, bottomtitle=3pt}
\newtcolorbox{rechtbox}{colback=lightgray, colframe=darkgray, fonttitle=\bfseries, title=Rechtsgrundlage / \textit{Legal Basis}, toptitle=3pt, bottomtitle=3pt}
\newtcolorbox{kritischbox}{colback=spanisch!15, colframe=spanisch, boxrule=3pt, fonttitle=\bfseries\large, title=RECHTLICH VERBINDLICH! / \textit{LEGALLY MANDATORY!}, coltitle=white, colbacktitle=spanisch, toptitle=5pt, bottomtitle=5pt}

\newcommand{\cmark}{\textcolor{success}{\ding{51}}}
\newcommand{\xmark}{\textcolor{spanisch}{\ding{55}}}

\begin{document}

%% TITLE PAGE
\begin{titlepage}
\begin{center}
\vspace*{1.5cm}

{\Huge\bfseries\textcolor{spaet}{SPANISCH}}\\[0.2cm]
{\Huge\bfseries\textcolor{spaet}{SPÄTBEGINNER}}\\[0.3cm]
{\large\color{englishgray}\textit{SPANISH LATE BEGINNERS}}\\[0.8cm]

{\LARGE\textcolor{darkgray}{Vollständige Übersicht für Mecklenburg-Vorpommern}}\\[0.2cm]
{\large\color{englishgray}\textit{Complete Guide for Mecklenburg-Western Pomerania}}\\[0.5cm]
{\large Klassen 10 -- 12 | Gesamtschule \& Gymnasium}\\[0.1cm]
{\small\color{englishgray}\textit{Grades 10--12 | Comprehensive School \& Grammar School}}\\[1.2cm]

\begin{tikzpicture}[scale=1.3]
\node[draw=spaet, fill=lightspaet, minimum width=2cm, minimum height=1.2cm, rounded corners=3pt, line width=1.5pt] (kl10) at (0,0) {\textbf{Klasse 10}\\{\scriptsize A1 $\rightarrow$ A2}};
\node[draw=spaet, fill=lightspaet, minimum width=2cm, minimum height=1.2cm, rounded corners=3pt, line width=1.5pt] (kl11) at (4,0) {\textbf{Klasse 11}\\{\scriptsize A2 $\rightarrow$ A2+}};
\node[draw=spaet, fill=spaet, text=white, minimum width=2cm, minimum height=1.2cm, rounded corners=3pt, line width=1.5pt] (kl12) at (8,0) {\textbf{Klasse 12}\\{\scriptsize A2+ $\rightarrow$ B1(+)}};

\draw[-{Stealth[length=4mm]}, line width=2pt, color=spaet] (kl10) -- (kl11);
\draw[-{Stealth[length=4mm]}, line width=2pt, color=spaet] (kl11) -- (kl12);

\node[below=0.6cm of kl10] {\footnotesize 4 WStd};
\node[below=0.6cm of kl11] {\footnotesize 4 WStd};
\node[below=0.6cm of kl12] {\footnotesize 4 WStd $\rightarrow$ Abitur};
\end{tikzpicture}

\vspace{1.5cm}

\fcolorbox{spaet}{lightspaet}{\parbox{0.9\textwidth}{\centering
\textbf{Dieses Dokument enthält: / \textit{This document contains:}}\\[4pt]
\small
\begin{tabular}{ll}
\textbullet\ Rechtlicher Rahmen / \textit{Legal Framework} & \textbullet\ Prüfungsformate / \textit{Exam Formats} \\
\textbullet\ Unterrichtsplanung / \textit{Lesson Planning} & \textbullet\ Lehrwerk-Analyse / \textit{Textbook Analysis} \\
\textbullet\ Themenfelder / \textit{Topic Areas} & \textbullet\ Handlungsempfehlungen / \textit{Recommendations} \\
\end{tabular}\\[6pt]
\textcolor{spanisch}{\textbf{WICHTIG: Bewertungsbögen sind rechtlich verbindlich!}}\\
{\scriptsize\color{englishgray}\textit{IMPORTANT: Assessment forms are legally mandatory!}}
}}

\vfill

{\small Version 2.0 | Januar 2026 | Bilingual DE/EN}\\
{\footnotesize Lernpfade | Spanisch Spätbeginner MV}\\
{\scriptsize\color{englishgray}\textit{Learning Paths | Spanish Late Beginners MV}}

\end{center}
\end{titlepage}

%% TABLE OF CONTENTS
\tableofcontents
\newpage

%% ============================================================================
%% TEIL 1: GRUNDLAGEN
%% ============================================================================

\section{Rechtlicher Rahmen und Definition}

\subsection{Was bedeutet ``Spätbeginner''?}

Der Begriff \textbf{``Spätbeginner''} ist eine \textit{informelle Bezeichnung}. Die offiziellen Dokumente des Landes Mecklenburg-Vorpommern verwenden:

\begin{wichtigbox}
\textbf{``In der Einführungsphase neu begonnene Fremdsprache''}\\[4pt]
Definition: Spanisch ab Klassenstufe 10 (Einführungsphase) \textbf{ohne Vorkenntnisse}.\\[4pt]
\textbf{Kritisch:} Es existiert \textbf{KEIN separater Rahmenplan} für Spätbeginner!\\
Der gleiche Rahmenplan wie für fortgeführte Fremdsprachen gilt.
\end{wichtigbox}

\subsection{Die maßgeblichen Rechtsquellen}

\begin{center}
\small
\begin{tabular}{|p{5cm}|p{3.5cm}|p{5cm}|}
\hline
\rowcolor{spaet}\color{white}\textbf{Dokument} & \color{white}\textbf{Herausgeber} & \color{white}\textbf{Relevanz für Spätbeginner} \\
\hline
Rahmenplan Spanisch Qualifikationsphase (2019) & Ministerium für Bildung M-V & Verbindliche Themenfelder, Kompetenzen \\
\hline
Oberstufen- und Abiturprüfungsverordnung (APVO M-V, 2019/2025) & Land M-V & Stundentafel (4 WStd!), Prüfungsformate, Klausurdauer \\
\hline
Handreichung Sprechprüfung (2020) & Institut für Qualitätsentwicklung & Format, Bewertungsbögen A2/B1/B2 \\
\hline
Handreichung Paarprüfung Abitur (2025) & Institut für Qualitätsentwicklung & Mündliche Abiturprüfung \\
\hline
VV Sprechprüfung (24.07.2020) & Mitteilungsblatt Nr. 5/2020 & Inkraftsetzung der Handreichung \\
\hline
VV Paarprüfung (03.12.2025) & Mitteilungsblatt Nr. 14/2025 & Inkraftsetzung Abitur-Paarprüfung \\
\hline
\end{tabular}
\end{center}

\subsection{Spätbeginner = Eigene Kategorie}

\begin{rechtbox}
\textbf{APVO M-V:} \textit{``Eine neu beginnende Fremdsprache wird \textbf{vierstündig} unterrichtet.''}
\end{rechtbox}

Die neu begonnene Fremdsprache ist \textbf{weder} Grundkurs (3 WStd) \textbf{noch} Leistungskurs (5 WStd), sondern eine \textbf{eigenständige Kategorie}:

\begin{center}
\begin{tabular}{|l|c|c|c|}
\hline
\rowcolor{lightgray}
\textbf{Kursart} & \textbf{Wochenstunden} & \textbf{LK möglich?} & \textbf{Zielniveau Abitur} \\
\hline
Grundkurs (fortgeführt) & 3 & Nein & B2 \\
\hline
Leistungskurs (fortgeführt) & 5 & Ja & B2+ \\
\hline
\rowcolor{lightspaet}
\textbf{Neu begonnene FS (Spätbeginner)} & \textbf{4} & \textbf{Nein} & \textbf{B1(+)} \\
\hline
\end{tabular}
\end{center}

\newpage

\section{GER-Progression und Stundentafel}

\subsection{Zielniveaus nach Klassenstufe}

\begin{center}
\begin{tikzpicture}[
    node distance=1.2cm,
    box/.style={draw=spaet, fill=lightspaet, minimum width=2.5cm, minimum height=1.3cm, align=center, rounded corners=3pt, line width=1pt},
    boxfinal/.style={draw=spaet, fill=spaet, text=white, minimum width=2.5cm, minimum height=1.3cm, align=center, rounded corners=3pt, line width=1pt},
    arrow/.style={-{Stealth[length=4mm]}, thick, color=spaet}
]
\node[box] (start) {\textbf{A1}\\{\scriptsize Beginn Kl. 10}};
\node[box, right=of start] (a2) {\textbf{A2}\\{\scriptsize Ende Kl. 10}};
\node[box, right=of a2] (a2p) {\textbf{A2+}\\{\scriptsize Ende Kl. 11}};
\node[boxfinal, right=of a2p] (b1) {\textbf{B1(+)}\\{\scriptsize Abitur}};

\draw[arrow] (start) -- (a2);
\draw[arrow] (a2) -- (a2p);
\draw[arrow] (a2p) -- (b1);

\node[below=0.6cm of start, align=center] {\scriptsize 160 Stunden\\{\tiny (4 WStd $\times$ 40 Wochen)}};
\node[below=0.6cm of a2, align=center] {\scriptsize 160 Stunden};
\node[below=0.6cm of a2p, align=center] {\scriptsize 160 Stunden};
\node[below=0.6cm of b1, align=center] {\scriptsize $\Sigma$ \textbf{480 Stunden}};
\end{tikzpicture}
\end{center}

\begin{rechtbox}
\textbf{Handreichung Sprechprüfung (2020), S. 3:}\\
\textit{``In der Einführungsphase neu begonnene Fremdsprachen orientieren sich am Ende der Einführungsphase am Niveau \textbf{A2} und am Ende der Qualifikationsphase am Zielniveau \textbf{B1(+)}.''}
\end{rechtbox}

\subsection{Vergleich: Spätbeginner vs. Fortgeführte Fremdsprache}

\begin{center}
\begin{tabular}{|l|c|c|c|c|}
\hline
\rowcolor{lightgray}
\textbf{Phase} & \multicolumn{2}{c|}{\textbf{Fortgeführte FS}} & \multicolumn{2}{c|}{\textbf{Spätbeginner}} \\
\cline{2-5}
\rowcolor{lightgray}
& Niveau & WStd & Niveau & WStd \\
\hline
Einführungsphase (Kl. 10) & B1 & 3 oder 5 & \textcolor{spaet}{\textbf{A2}} & \textcolor{spaet}{\textbf{4}} \\
\hline
Qualifikationsphase 1 (Kl. 11) & B1+ & 3 oder 5 & \textcolor{spaet}{\textbf{A2+}} & \textcolor{spaet}{\textbf{4}} \\
\hline
Qualifikationsphase 2 (Kl. 12) & B2 & 3 oder 5 & \textcolor{spaet}{\textbf{B1(+)}} & \textcolor{spaet}{\textbf{4}} \\
\hline
\textbf{Abitur} & \textbf{B2} & & \textcolor{spaet}{\textbf{B1(+)}} & \\
\hline
\end{tabular}
\end{center}

\newpage

%% ============================================================================
%% TEIL 2: UNTERRICHT NACH KLASSENSTUFE
%% ============================================================================

\section{Klasse 10 (Einführungsphase)}

\subsection{Überblick}

\begin{center}
\begin{tabular}{|p{4cm}|p{9.5cm}|}
\hline
\rowcolor{spaet}\color{white}\textbf{Merkmal} & \color{white}\textbf{Details} \\
\hline
\textbf{Zielniveau} & A2 (Gemeinsamer Europäischer Referenzrahmen) \\
\hline
\textbf{Wochenstunden} & 4 \\
\hline
\textbf{Jahresstunden} & ca. 160 (4 $\times$ 40 Wochen) \\
\hline
\textbf{Fokus} & Grundlagen aufbauen, kommunikative Basiskompetenz entwickeln \\
\hline
\textbf{Klausuren} & 2 pro Jahr, mindestens 45 Minuten (90 Min. bei Aufsätzen) \\
\hline
\textbf{Sprechprüfung} & Optional (empfohlen: noch nicht in Klasse 10) \\
\hline
\textbf{Lehrwerk} & A\_tope.com: Unidades 1--4 \\
\hline
\end{tabular}
\end{center}

\subsection{Halbjahresplanung}

\subsubsection{Halbjahr 10.1 (August--Januar)}

\begin{center}
\small
\begin{tabular}{|c|p{5cm}|p{6cm}|}
\hline
\rowcolor{lightgray}
\textbf{Wochen} & \textbf{Thema} & \textbf{A\_tope.com} \\
\hline
1--4 & Grundlagen (Aussprache, Alphabet, Basisgrammatik) & Hablamos español + Unidad 1 (S. 8--21) \\
\hline
5--8 & Yo y mi entorno (sich vorstellen, Familie) & Unidad 2 (S. 22--39) \\
\hline
9--12 & La vida cotidiana (Alltag, Uhrzeit, Aktivitäten) & Unidad 3 (S. 40--51) \\
\hline
13--16 & Mi colegio y mis amigos + \textbf{Klausur 1} & Unidad 4 (S. 52--69) \\
\hline
\end{tabular}
\end{center}

\subsubsection{Halbjahr 10.2 (Februar--Juli)}

\begin{center}
\small
\begin{tabular}{|c|p{5cm}|p{6cm}|}
\hline
\rowcolor{lightgray}
\textbf{Wochen} & \textbf{Thema} & \textbf{A\_tope.com} \\
\hline
1--4 & España: geografía básica & Unidad 5 (S. 70--81) \\
\hline
5--8 & Tradiciones y fiestas & Panorama 1 (S. 68--69) \\
\hline
9--12 & De compras y en el restaurante & Unidad 5 (S. 79--81) \\
\hline
13--16 & Wiederholung + \textbf{Klausur 2} & Evaluación 1 \\
\hline
\end{tabular}
\end{center}

\subsection{Grammatik Klasse 10}

\begin{center}
\begin{tabular}{|l|p{9cm}|}
\hline
\rowcolor{lightgray}
\textbf{Struktur} & \textbf{Beschreibung} \\
\hline
Presente de indicativo & Regelmäßige Verben + wichtige unregelmäßige (ser, estar, tener, ir, hacer) \\
\hline
Artículos, género, número & Bestimmte/unbestimmte Artikel, Genus und Numerus \\
\hline
Possessivbegleiter & mi, tu, su, nuestro, vuestro, su \\
\hline
Pretérito perfecto & Einführung (Perfekt mit haber + Partizip) \\
\hline
Futuro próximo & ir a + Infinitiv \\
\hline
Präpositionen & Grundlegende Ortsangaben (en, a, de, con, para, por) \\
\hline
gustar & gustar + Infinitiv, gustar + Substantiv \\
\hline
\end{tabular}
\end{center}

\begin{empfehlungbox}
\textbf{Klasse 10:} A\_tope.com deckt die Grundlagen zu \textbf{ca. 90\%} ab. Das Lehrwerk ist für diese Klassenstufe \textbf{ausreichend} ohne Ergänzungen.
\end{empfehlungbox}

\newpage

\section{Klasse 11 (Qualifikationsphase 1)}

\subsection{Überblick}

\begin{center}
\begin{tabular}{|p{4cm}|p{9.5cm}|}
\hline
\rowcolor{spaet}\color{white}\textbf{Merkmal} & \color{white}\textbf{Details} \\
\hline
\textbf{Zielniveau} & A2+ $\rightarrow$ B1 (Progression) \\
\hline
\textbf{Wochenstunden} & 4 \\
\hline
\textbf{Jahresstunden} & ca. 160 \\
\hline
\textbf{Themenfelder} & 1 (El individuo) + 2 (Identidad nacional) -- vereinfacht \\
\hline
\textbf{Klausuren} & 2 pro Jahr, mindestens 90 Minuten \\
\hline
\textbf{Sprechprüfung} & \textbf{EMPFOHLEN} (ersetzt 1 Klausur) \\
\hline
\textbf{Lehrwerk} & A\_tope.com: Unidades 5--8 + Module \\
\hline
\end{tabular}
\end{center}

\subsection{Halbjahresplanung}

\subsubsection{Halbjahr 11.1: Themenfeld 1 -- El individuo y la sociedad}

\begin{center}
\small
\begin{tabular}{|c|p{5cm}|p{6cm}|}
\hline
\rowcolor{lightgray}
\textbf{Wochen} & \textbf{Thema} & \textbf{A\_tope.com} \\
\hline
1--5 & Perú: ayer y hoy (indefinido) & Unidad 6 (S. 82--99) \\
\hline
6--9 & Buscando trabajo (Berufe, Zukunft) & Unidad 7 (S. 100--109) \\
\hline
10--12 & Relaciones y emociones & Módulo: Una decisión importante \\
\hline
13--16 & \textbf{Klausur 1} + Wiederholung & Evaluación 2 \\
\hline
\end{tabular}
\end{center}

\subsubsection{Halbjahr 11.2: Themenfeld 2 -- La identidad nacional}

\begin{center}
\small
\begin{tabular}{|c|p{5cm}|p{6cm}|}
\hline
\rowcolor{lightgray}
\textbf{Wochen} & \textbf{Thema} & \textbf{A\_tope.com} \\
\hline
1--5 & Andalucía (imperfecto) & Unidad 8 (S. 110--125) \\
\hline
6--9 & España: regiones y cultura & Panorama 3 (Comunidades Autónomas) \\
\hline
10--12 & Diversidad cultural & Módulo: Vivir la diversidad \\
\hline
13--16 & \textbf{Sprechprüfung} (B1) oder Klausur 2 & --- \\
\hline
\end{tabular}
\end{center}

\subsection{Grammatik Klasse 11}

\begin{center}
\begin{tabular}{|l|p{9cm}|}
\hline
\rowcolor{lightgray}
\textbf{Struktur} & \textbf{Beschreibung} \\
\hline
Pretérito indefinido & Vollständige Einführung (regelmäßig + unregelmäßig) \\
\hline
Pretérito imperfecto & Formen und Verwendung \\
\hline
Kontrast indefinido/imperfecto & Unterscheidung in Erzählungen \\
\hline
Subjuntivo presente & Einführung (que + Subjuntivo nach Wunsch, Aufforderung) \\
\hline
Condicional simple & Konditional für höfliche Bitten, Wünsche \\
\hline
Relativsätze & que, donde, quien \\
\hline
\end{tabular}
\end{center}

\begin{empfehlungbox}
\textbf{Sprechprüfung in Klasse 11:} Die Handreichung empfiehlt, die Sprechprüfung im \textbf{2. Halbjahr Klasse 11} durchzuführen (B1-Bogen). So bleibt in Klasse 12 mehr Zeit für die Abitur-Vorbereitung und die Korrekturentlastung wird erreicht.
\end{empfehlungbox}

\newpage

\section{Klasse 12 (Qualifikationsphase 2 und Abitur)}

\subsection{Überblick}

\begin{center}
\begin{tabular}{|p{4cm}|p{9.5cm}|}
\hline
\rowcolor{spaet}\color{white}\textbf{Merkmal} & \color{white}\textbf{Details} \\
\hline
\textbf{Zielniveau} & B1(+) (Abiturniveau) \\
\hline
\textbf{Wochenstunden} & 4 \\
\hline
\textbf{Jahresstunden} & ca. 160 \\
\hline
\textbf{Themenfelder} & 3 (Aspectos actuales) + 4 (Desafíos modernos) \\
\hline
\textbf{Klausuren} & 2 pro Jahr, mindestens 90 Minuten (letzte: abiturnah) \\
\hline
\textbf{Sprechprüfung} & \textbf{PFLICHT}, falls noch nicht erfolgt \\
\hline
\textbf{Lehrwerk} & A\_tope.com: Unidad 8 + Module -- \textcolor{spanisch}{\textbf{ERGÄNZUNGEN ERFORDERLICH!}} \\
\hline
\end{tabular}
\end{center}

\subsection{Halbjahresplanung}

\begin{tcolorbox}[colback=lightred, colframe=spanisch, title=\textbf{WICHTIG}]
\textbf{KRITISCH:} Das Lehrwerk A\_tope.com deckt die Themenfelder 3 und 4 \textbf{UNZUREICHEND} ab!

\textbullet\ \textbf{Migration:} 0\% Deckung -- nicht vorhanden!\\
\textbullet\ \textbf{Umwelt:} ca. 15\% -- nur ``turismo sostenible'' in Unidad 8\\
\textbullet\ \textbf{Digitalisierung:} nur Vokabular, keine kritische Reflexion

\textbf{$\Rightarrow$ Ergänzungsmaterial ist ZWINGEND ERFORDERLICH!}
\end{tcolorbox}

\subsubsection{Halbjahr 12.1: Themenfeld 3 -- Aspectos actuales de la política y la sociedad}

\begin{center}
\small
\begin{tabular}{|c|p{4.5cm}|p{6cm}|}
\hline
\rowcolor{lightgray}
\textbf{Wochen} & \textbf{Thema} & \textbf{Material} \\
\hline
1--5 & \textbf{Migración y sus causas} & \textcolor{spanisch}{\textbf{Lektüre \textit{Abdel}} (Enrique Páez) -- ERGÄNZUNG!} \\
\hline
6--9 & Globalización & Módulo: El comercio justo (S. 128--129) \\
\hline
10--12 & \textbf{Klausur 1} (abiturnah) & --- \\
\hline
13--16 & Wiederholung Themenfeld 1+2 & --- \\
\hline
\end{tabular}
\end{center}

\subsubsection{Halbjahr 12.2: Themenfeld 4 -- Desafíos modernos + Abitur}

\begin{center}
\small
\begin{tabular}{|c|p{4.5cm}|p{6cm}|}
\hline
\rowcolor{lightgray}
\textbf{Wochen} & \textbf{Thema} & \textbf{Material} \\
\hline
1--3 & \textbf{Medio ambiente} & \textcolor{spanisch}{\textbf{Aktualtexte -- ERGÄNZUNG!}} \\
\hline
4--5 & Tecnología y redes sociales & Módulo: ¿Te comunicas? + Aktualtexte \\
\hline
6--10 & \textbf{Abitur-Vorbereitung} & Alle 4 Themenfelder wiederholen \\
\hline
11+ & \textbf{Abiturprüfung} & --- \\
\hline
\end{tabular}
\end{center}

\subsection{Erforderliche Ergänzungen für Klasse 12}

\begin{center}
\begin{tabular}{|p{3.5cm}|p{5cm}|p{4.5cm}|}
\hline
\rowcolor{spanisch}\color{white}\textbf{Verbindlicher Inhalt} & \color{white}\textbf{Empfohlene Ergänzung} & \color{white}\textbf{Quelle} \\
\hline
\textbf{Movimientos migratorios} (Migration) & Lektüre \textit{Abdel} (Enrique Páez) & Literaturempfehlungen M-V \\
\hline
\textbf{Tendencias ecológicas} (Umwelt) & Aktualtexte (cambio climático, contaminación) & Rahmenplan S. 19 \\
\hline
\textbf{Desafíos de la era digital} & Aktualtexte zur kritischen Medienreflexion & Rahmenplan S. 19 \\
\hline
Globalización & A\_tope Módulo ``El comercio justo'' + Aktualtexte & --- \\
\hline
\end{tabular}
\end{center}

\begin{empfehlungbox}
\textbf{Top-Empfehlung: Lektüre \textit{Abdel} (Enrique Páez)}\\[4pt]
16-jähriger Protagonist aus Marokko. Thema: Illegale Immigration nach Spanien.\\
Sprachlich zugänglich (A2/B1), chronologische Erzählung, ermöglicht Identifikation.\\[4pt]
\textbf{Praktisch unverzichtbar} für Themenfeld 3, da A\_tope.com Migration nicht behandelt!
\end{empfehlungbox}

\newpage

%% ============================================================================
%% TEIL 3: THEMENFELDER
%% ============================================================================

\section{Die vier verbindlichen Themenfelder}

Der Rahmenplan Spanisch Qualifikationsphase M-V (2019) definiert \textbf{vier verbindliche Themenfelder}, die für alle Kurse gelten -- auch für Spätbeginner. Die Tiefe wird dem Sprachniveau angepasst.

\begin{center}
\begin{tabular}{|c|p{5.5cm}|c|c|c|}
\hline
\rowcolor{spaet}\color{white}\textbf{Nr.} & \color{white}\textbf{Themenfeld} & \color{white}\textbf{Fortgef.} & \color{white}\textbf{Spätbeg.} & \color{white}\textbf{A\_tope} \\
\hline
1 & El individuo y la sociedad & 45h & 30--35h & \textcolor{success}{\textbf{90\%}} \\
\hline
2 & La identidad nacional y la diversidad cultural & 45h & 30--35h & \textcolor{warning}{\textbf{70\%}} \\
\hline
3 & Aspectos actuales de la política y la sociedad & 30h & 20--25h & \textcolor{spanisch}{\textbf{0\%}} \\
\hline
4 & Desafíos modernos & 30h & 20--25h & \textcolor{spanisch}{\textbf{15\%}} \\
\hline
\rowcolor{lightspaet}
& \textbf{GESAMT} & \textbf{150h} & \textbf{100--120h} & \\
\hline
\end{tabular}
\end{center}

\subsection{Verbindlich vs. Empfohlen}

\begin{rechtbox}
\textbf{Rahmenplan, S. 1:}\\
\textit{``Die [...] benannten Themen und Ziele füllen ca. \textbf{80\%} der zur Verfügung stehenden Unterrichtszeit. [...] 20\% stehen für eigene Unterrichtsgestaltung zur Verfügung.''}
\end{rechtbox}

\begin{center}
\begin{tabular}{|p{6cm}|p{7cm}|}
\hline
\rowcolor{lightgray}
\textbf{Linke Spalte im RP} & \textbf{Rechte Spalte im RP} \\
\hline
\textbf{VERBINDLICH} & EMPFOHLEN \\
\hline
Muss behandelt werden & Hinweise und Anregungen \\
\hline
Prüfungsrelevant & Beispiele zur Schwerpunktsetzung \\
\hline
\end{tabular}
\end{center}

\newpage

%% ============================================================================
%% TEIL 4: PRÜFUNGSFORMATE
%% ============================================================================

\section{Prüfungsformate}

\subsection{Klausuren}

\begin{center}
\begin{tabular}{|l|c|c|c|c|}
\hline
\rowcolor{spaet}\color{white}\textbf{Phase} & \color{white}\textbf{Anzahl/Jahr} & \color{white}\textbf{Mindestdauer} & \color{white}\textbf{Niveau} & \color{white}\textbf{Ersetzbar?} \\
\hline
Einführungsphase (Kl. 10) & 2 & 45 Min. (90 bei Aufsätzen) & A2 & 1$\times$/Jahr \\
\hline
Qualifikationsphase (Kl. 11) & 2 & 90 Min. & A2+/B1 & 1$\times$/Jahr \\
\hline
Qualifikationsphase (Kl. 12) & 2 (1 im Abitur-Semester) & 90 Min. & B1 & 1$\times$/Jahr \\
\hline
\end{tabular}
\end{center}

\begin{wichtigbox}
\textbf{§ 22 Abs. 2 APVO:} Die \textbf{Klausur unter abiturähnlichen Bedingungen (Q4)} kann \textbf{NICHT} durch eine Sprechprüfung ersetzt werden!
\end{wichtigbox}

\subsection{Sprechprüfung (Komplexe Leistungsermittlung Sprechen)}

\subsubsection{Rechtliche Grundlage}

\begin{rechtbox}
\textbf{§ 22 Abs. 5 APVO M-V:}\\
\textit{``In den modernen Fremdsprachen erfolgt für jede Schülerin und jeden Schüler \textbf{mindestens eine} komplexe Leistungsermittlung im Kompetenzbereich Sprechen.''}
\end{rechtbox}

\subsubsection{Format: Drei Teile (Tandem-Prüfung)}

\begin{center}
\begin{tabular}{|c|p{3.5cm}|c|p{6cm}|}
\hline
\rowcolor{spaet}\color{white}\textbf{Teil} & \color{white}\textbf{Name} & \color{white}\textbf{Dauer} & \color{white}\textbf{Beschreibung} \\
\hline
1 & Interview & 2 Min. & Aufwärmen, persönliche Fragen -- \textbf{KEINE Bewertung!} \\
\hline
2 & Monolog & 4 Min./SuS & Impulsgesteuerte Präsentation (Bild, Text, Grafik) \\
\hline
3 & Dialog & 5 Min./SuS & Diskussion oder Rollenspiel zu einem Thema \\
\hline
\multicolumn{2}{|c|}{\textbf{Gesamt}} & \textbf{$\approx$20 Min.} & Tandem-Prüfung (2 SuS gleichzeitig) \\
\hline
\end{tabular}
\end{center}

\subsubsection{Bewertung}

\begin{center}
\begin{tabular}{|p{6cm}|c|}
\hline
\rowcolor{lightgray}
\textbf{Kriterium} & \textbf{Gewichtung} \\
\hline
\textbf{Sprachliche Leistung} & \textbf{60\%} \\
\quad Aussprache und Intonation & \\
\quad Wortschatz und Ausdruck & \\
\quad Grammatische Korrektheit & \\
\quad Flüssigkeit und Interaktion & \\
\hline
\textbf{Inhaltliche Leistung} & \textbf{40\%} \\
\quad Aufgabenerfüllung & \\
\quad Argumentation & \\
\quad Themenentwicklung & \\
\hline
\end{tabular}
\end{center}

\subsubsection{Bewertungsbögen nach Niveau (für Spätbeginner)}
{\small\color{englishgray}\textit{Assessment Forms by Level (for Late Beginners)}}

\begin{kritischbox}
\textbf{ACHTUNG: Die Excel-Bewertungsbögen (Anlage II)\footnote{Anlage II = Anhang II der Handreichung zur KLE Sprechen / \textit{Appendix II of the KLE Speaking Guidelines}} sind RECHTLICH VERBINDLICH!}\\[0.2cm]
{\small\color{englishgray}\textit{WARNING: The Excel assessment forms (Appendix II) are LEGALLY MANDATORY!}}\\[0.4cm]

\textbf{Rechtsgrundlage / \textit{Legal Basis}:}\\
Handreichung zur KLE\footnote{KLE = Komplexe Leistungsermittlung / \textit{Complex Performance Assessment}} Sprechen, S. 13:\\
\glqq Die \textbf{verbindlich zu nutzenden} Bewertungsbögen für Gesprächsleitungen und Protokollanten (Anlage II) sind im SIP-Datentauschportal\footnote{SIP = Schulinformations- und Planungssystem M-V / \textit{School Information and Planning System M-V}} unter `Sprechleistung\_Moderne Fremdsprachen' abrufbar.\grqq\\[0.3cm]

\textbf{Konsequenz bei Nichtverwendung / \textit{Consequences of Non-Use}:}\\
$\bullet$ Prüfung ist \textbf{formal anfechtbar} / \textit{Examination is formally contestable}\\
$\bullet$ Noten können bei Widerspruch \textbf{aufgehoben} werden / \textit{Grades can be annulled upon appeal}\\
$\bullet$ \textbf{Eigene Bewertungsbögen sind NICHT zulässig!} / \textit{Custom forms are NOT permitted!}
\end{kritischbox}

\vspace{0.4cm}

\begin{center}
\begin{tabular}{|l|c|p{4cm}|}
\hline
\rowcolor{spaet}\color{white}\textbf{Phase} & \color{white}\textbf{Bogen} & \color{white}\textbf{GER\footnote{GER = Gemeinsamer Europäischer Referenzrahmen / \textit{CEFR = Common European Framework of Reference}}-Niveau} \\
\hline
Einführungsphase (Kl. 10) & \textcolor{spaet}{\textbf{A2}} & Elementare Sprachverwendung \\
{\scriptsize\color{englishgray}\textit{Introductory Phase (Gr. 10)}} & & {\scriptsize\color{englishgray}\textit{Basic User}} \\
\hline
Qualifikationsphase (Kl. 11/12) & \textcolor{spaet}{\textbf{B1}} & Selbstständige Sprachverwendung \\
{\scriptsize\color{englishgray}\textit{Qualification Phase (Gr. 11/12)}} & & {\scriptsize\color{englishgray}\textit{Independent User}} \\
\hline
\multicolumn{3}{|l|}{\textit{Fortgeführte FS\footnote{FS = Fremdsprache / \textit{FL = Foreign Language}}: B2-Bogen / \textit{Continuing FL: B2 form}}} \\
\hline
\end{tabular}
\end{center}

\vspace{0.4cm}

\textbf{\large Wie erhalte ich die Bewertungsbögen? / \textit{How do I obtain the assessment forms?}}

\vspace{0.3cm}

\begin{center}
\small
\begin{tabular}{|c|p{3.5cm}|p{7.5cm}|}
\hline
\rowcolor{spaet}\color{white}\textbf{Nr.} & \color{white}\textbf{Beschaffungsweg} & \color{white}\textbf{Details / \textit{Details}} \\
\hline
\textbf{1} & \textbf{Schulleitung}\newline {\scriptsize\color{englishgray}\textit{School Administration}} & Oberstufenkoordinator/in\footnote{OStKo = Oberstufenkoordinator / \textit{Upper School Coordinator}} oder Schulleiter/in anfragen -- haben Zugang zu allen Portalen\newline {\scriptsize\color{englishgray}\textit{Ask coordinator or principal -- they have access to all portals}} \\
\hline
\textbf{2} & \textbf{Fachkonferenzleitung}\newline {\scriptsize\color{englishgray}\textit{Subject Conference}} & FK\footnote{FK = Fachkonferenz / \textit{SC = Subject Conference}}-Leitung Fremdsprachen sollte Materialien für alle Kolleg/innen haben\newline {\scriptsize\color{englishgray}\textit{FL subject leader should have materials for all colleagues}} \\
\hline
\textbf{3} & \textbf{SIP-Portal}\newline {\scriptsize\color{englishgray}\textit{SIP Portal}} & Zugang über Schulserver, Pfad: \texttt{Sprechleistung\_Moderne Fremdsprachen}\newline {\scriptsize\color{englishgray}\textit{Access via school server}} \\
\hline
\textbf{4} & \textbf{itslearning-Kurs}\newline {\scriptsize\color{englishgray}\textit{itslearning Course}} & Kurs: \glqq Sprechen testen in den Fremdsprachen SEK II\grqq\newline {\scriptsize\color{englishgray}\textit{Course: ``Testing Speaking in FL Upper Secondary''}} \\
\hline
\textbf{5} & \textbf{IQ M-V\footnote{IQ M-V = Institut für Qualitätsentwicklung M-V / \textit{Institute for Quality Development M-V}} direkt}\newline {\scriptsize\color{englishgray}\textit{IQ M-V directly}} & E-Mail: \texttt{iq@bm.mv-regierung.de} -- freundlich anfragen!\newline {\scriptsize\color{englishgray}\textit{Email request -- usually very helpful!}} \\
\hline
\textbf{6} & \textbf{Fachberater Spanisch}\newline {\scriptsize\color{englishgray}\textit{Subject Advisor}} & Über Schulamt\footnote{Schulamt = Staatliches Schulamt / \textit{State School Authority}} erfragen -- kennt alle Bezugsquellen\newline {\scriptsize\color{englishgray}\textit{Inquire via school authority -- knows all sources}} \\
\hline
\end{tabular}
\end{center}

\begin{empfehlungbox}
\textbf{EMPFOHLENE REIHENFOLGE / \textit{RECOMMENDED ORDER}:}\\[0.2cm]
\textbf{1.} Schulleitung / FK-Leitung fragen (schnellster Weg)\\
{\small\color{englishgray}\textit{Ask school admin / subject leader (fastest way)}}\\[0.1cm]
\textbf{2.} Falls erfolglos: IQ M-V direkt anschreiben\\
{\small\color{englishgray}\textit{If unsuccessful: contact IQ M-V directly}}\\[0.3cm]
\textbf{Weitere Materialien / \textit{Additional Materials}:}\\
$\bullet$ Regieanweisungen, Beispielplanungen: itslearning-Kurs SEK II\\
$\bullet$ Abitur-Paarprüfung\footnote{Paarprüfung = mündliche Prüfung mit zwei Prüflingen / \textit{Pair Examination = oral exam with two candidates}} (ab 2025/26): itslearning-Kurs \glqq Paarprüfung im Abitur\grqq\\
$\bullet$ Rahmenpläne: \url{https://www.bildung-mv.de}
\end{empfehlungbox}

\newpage

\subsection{Abiturprüfung}

\subsubsection{Schriftliches Abitur}

\begin{center}
\begin{tabular}{|p{5cm}|p{8.5cm}|}
\hline
\rowcolor{spaet}\color{white}\textbf{Merkmal} & \color{white}\textbf{Details für Spätbeginner} \\
\hline
\textbf{Dauer} & 180 Minuten (nur grundlegendes Niveau, kein LK!) \\
\hline
\textbf{Niveau} & B1(+) \\
\hline
\textbf{Prüfungsfach} & Wählbar als P4 (Kombination schriftlich+mündlich) \\
\hline
\textbf{Aufgabenformate} & Textaufgabe (Pflicht), Leseverstehen, Hörverstehen (opt.), Sprachmittlung (opt.) \\
\hline
\end{tabular}
\end{center}

\subsubsection{Anforderungsbereiche (Spätbeginner vs. Fortgeführt)}

\begin{center}
\begin{tabular}{|l|c|c|p{4cm}|}
\hline
\rowcolor{lightgray}
\textbf{Anforderungsbereich} & \textbf{Spätbeginner} & \textbf{Fortgeführt} & \textbf{Operatoren} \\
\hline
I (Reproduktion) & \textcolor{spaet}{\textbf{25\%}} & 20\% & nenne, beschreibe, gib an \\
\hline
II (Reorganisation) & \textcolor{spaet}{\textbf{50\%}} & 45\% & erkläre, vergleiche, analysiere \\
\hline
III (Transfer) & \textcolor{spaet}{\textbf{25\%}} & 35\% & beurteile, diskutiere, entwickle \\
\hline
\end{tabular}
\end{center}

\subsubsection{Mündliche Abiturprüfung (Paarprüfung)}

\textbf{Rechtliche Grundlage:} § 25 Abs. 8 (P4) und § 38 Abs. 6 (P5) APVO M-V, VV Paarprüfung vom 03.12.2025

\begin{center}
\begin{tabular}{|c|p{3.5cm}|c|p{6cm}|}
\hline
\rowcolor{spaet}\color{white}\textbf{Teil} & \color{white}\textbf{Name} & \color{white}\textbf{Dauer} & \color{white}\textbf{Beschreibung} \\
\hline
1 & Interview & 2 Min. & Kurzer Einstieg -- \textbf{KEINE Bewertung!} \\
\hline
2 & Monolog & 4 Min./Prüfling & Längerer Gesprächsbeitrag auf Basis eines Impulses \\
\hline
3 & Dialog & 5 Min./Prüfling & Interaktion zwischen beiden Prüflingen \\
\hline
\multicolumn{2}{|c|}{\textbf{Gesamt}} & \textbf{$\approx$30 Min.} & Vorbereitungszeit: 20 Min. (beide gemeinsam) \\
\hline
\end{tabular}
\end{center}

\newpage

%% ============================================================================
%% TEIL 5: ZUSAMMENFASSUNG
%% ============================================================================

\section{Zusammenfassung: Prüfungen für Spätbeginner}

\begin{center}
\begin{tabular}{|p{3.5cm}|c|c|c|p{4cm}|}
\hline
\rowcolor{spaet}\color{white}\textbf{Prüfung} & \color{white}\textbf{Anzahl} & \color{white}\textbf{Dauer} & \color{white}\textbf{Niveau} & \color{white}\textbf{Besonderheit} \\
\hline
Klausur Kl. 10 & 2/Jahr & 45--90 Min. & A2 & 1$\times$/Jahr ersetzbar \\
\hline
Klausur Kl. 11--12 & 2/Jahr & 90 Min. & A2+/B1 & 1$\times$/Jahr ersetzbar \\
\hline
\textbf{Sprechprüfung} & \textbf{$\geq$1 ges.} & 20 Min. & A2/B1 & \textbf{PFLICHT!} \\
\hline
Abitur schriftlich & 1 & 180 Min. & B1(+) & Nur grundlegendes Niveau \\
\hline
Abitur mündlich P5 & 1 & 30 Min. & B1 & Paarprüfung möglich \\
\hline
Abitur Kombi P4 & 1 & variabel & B1 & Verkürzt schriftl. + mündl. \\
\hline
\end{tabular}
\end{center}

\section{Was ist verbindlich, was ist erlaubt?}

\begin{center}
\begin{tabular}{|p{5.5cm}|c|p{5.5cm}|}
\hline
\rowcolor{spaet}\color{white}\textbf{Aspekt} & \color{white}\textbf{Status} & \color{white}\textbf{Quelle} \\
\hline
4 Wochenstunden für Spätbeginner & \textbf{Verbindlich} & APVO M-V \\
\hline
4 Themenfelder behandeln & \textbf{Verbindlich} & Rahmenplan S. 16--19 \\
\hline
Linke Spalte (verbindliche Inhalte) & \textbf{Verbindlich} & Rahmenplan S. 1 \\
\hline
Rechte Spalte (Hinweise) & Empfohlen & Rahmenplan S. 1 \\
\hline
Mindestens 1 Sprechprüfung & \textbf{Verbindlich} & § 22 Abs. 5 APVO \\
\hline
Reihenfolge der Themenfelder & Frei wählbar & Rahmenplan S. 1 \\
\hline
20\% der Unterrichtszeit & Frei gestaltbar & Rahmenplan S. 1 \\
\hline
Literatur (z.B. \textit{Abdel}) & Empfohlen & Literaturempfehlungen \\
\hline
Jahrgangsübergreifender Unterricht & \textbf{Erlaubt} & Schulgesetz M-V \\
\hline
\end{tabular}
\end{center}

\newpage

%% ============================================================================
%% TEIL 6: HANDLUNGSEMPFEHLUNGEN
%% ============================================================================

\section{Praktische Handlungsempfehlungen}

\subsection{Jahrgangsübergreifende Kurse}

An Gesamtschulen werden Spätbeginner-Kurse oft \textbf{jahrgangsübergreifend} unterrichtet (Kl. 10, 11, 12 gemeinsam).

\begin{tcolorbox}[colback=lightspaet, colframe=spaet, title=\textbf{HANDLUNGSEMPFEHLUNG}]
\textbf{Jahrgangsübergreifender Unterricht ist ZULÄSSIG}, wenn:

1. Alle SuS erhalten \textbf{4 Wochenstunden}\\
2. \textbf{Klausuren differenziert} nach Jahrgang (Kl. 10: A2, Kl. 11-12: B1)\\
3. Bewertung nach \textbf{individuellem Stand}\\
4. Mindestens 1 Sprechprüfung pro Schüler

Die Schulkonferenz muss dies beschließen.
\end{tcolorbox}

\subsection{Empfohlene Zeitleiste für Sprechprüfung}

\begin{center}
\begin{tikzpicture}[scale=0.9]
\draw[thick, ->, color=spaet] (0,0) -- (12,0);
\node at (0,-0.5) {\scriptsize Kl. 10};
\node at (4,-0.5) {\scriptsize Kl. 11};
\node at (8,-0.5) {\scriptsize Kl. 12};
\node at (12,-0.5) {\scriptsize Abitur};

\draw[thick, color=darkgray] (0,0.2) -- (0,-0.2);
\draw[thick, color=darkgray] (4,0.2) -- (4,-0.2);
\draw[thick, color=darkgray] (8,0.2) -- (8,-0.2);
\draw[thick, color=darkgray] (12,0.2) -- (12,-0.2);

\node[draw=spaet, fill=lightspaet, rounded corners=2pt, minimum height=0.8cm] at (6,1) {\textbf{Sprechprüfung (B1)}};
\draw[->, thick, color=spaet] (6,0.5) -- (6,0.1);
\node[below] at (6,-1) {\footnotesize \textbf{EMPFOHLEN:} 2. Halbjahr Kl. 11};
\end{tikzpicture}
\end{center}

\textbf{Begründung:} Im 4. Halbjahr korrigieren Englisch-Lehrkräfte die schriftlichen Abiturprüfungen. Durch Verlagerung der Spanisch-Sprechprüfung ins 2. Halbjahr Kl. 11 wird Korrekturentlastung erreicht.

\subsection{Checkliste: Materialplanung Klasse 12}

\begin{center}
\begin{tabular}{|c|p{8cm}|c|}
\hline
\rowcolor{lightgray}
& \textbf{Material} & \textbf{Vorhanden?} \\
\hline
$\square$ & Lektüre \textit{Abdel} (Enrique Páez) für Thema Migration & \\
\hline
$\square$ & Aktualtexte Umwelt (cambio climático, contaminación) & \\
\hline
$\square$ & Aktualtexte Digitalisierung (Chancen und Risiken) & \\
\hline
$\square$ & Wiederholungsmaterial Themenfelder 1+2 & \\
\hline
$\square$ & Abitur-Übungsklausuren auf B1-Niveau & \\
\hline
\end{tabular}
\end{center}

\subsection{Top-5 Empfehlungen}

\begin{enumerate}
\item \textbf{Sprechprüfung in Klasse 11} durchführen (nicht warten bis Klasse 12)
\item \textbf{Lektüre \textit{Abdel}} beschaffen -- unverzichtbar für Themenfeld 3
\item \textbf{Aktualtexte} für Themenfelder 3+4 sammeln (Rahmenplan-Lücken schließen)
\item Bei \textbf{jahrgangsübergreifenden Kursen}: Klausuren nach Niveau differenzieren
\item \textbf{Grammatik B2 (Pluscuamperfecto, Passiv)} ist für B1(+) NICHT zwingend erforderlich
\end{enumerate}

\vfill

\begin{center}
\fcolorbox{spaet}{lightspaet}{\parbox{0.9\textwidth}{\centering
\small
\textbf{Spanisch Spätbeginner -- Vollständige Übersicht MV}\\[4pt]
Version 1.0 | Januar 2026\\[4pt]
\textcolor{darkgray}{Lernpfade | Für Unterrichtsplanung und Prüfungsvorbereitung}
}}
\end{center}

\end{document}
